% =====================================================================
\chapter{Estimación por máxima verosimilitud y suavizado}
\label{tema:3}
% =====================================================================

Los parámetros de un modelo de $n$-gramas se estiman a partir de un corpus. Este tema demuestra que el estimador de máxima verosimilitud es la frecuencia relativa, analiza su principal defecto (las probabilidades nulas) y presenta dos remedios: el suavizado aditivo y la interpolación. Se sigue~\cite[§3.1.2 y §3.6]{jm3} y, para la interpretación bayesiana, \cite[§2.2]{bishop}.

\section{Verosimilitud de un corpus}

Sea un corpus formado por secuencias, cada una delimitada por \fin\ al principio y al final, y sea $N(u,v)$ el número de veces que el token $v$ sigue inmediatamente al token $u$. Se escribe $N(u)=\sum_{v}N(u,v)$ para el número de veces que $u$ aparece como token anterior.

\begin{definicion}[Verosimilitud]
Dada una matriz de transición $P$, la \emph{verosimilitud} del corpus es el producto de las probabilidades de todos sus bigramas,
\[
\mathcal{L}(P)=\prod_{u,v\in\V}P_{uv}^{\,N(u,v)},
\qquad
\ell(P)=\ln\mathcal{L}(P)=\sum_{u,v\in\V}N(u,v)\ln P_{uv},
\]
y $\ell$ se llama \emph{log-verosimilitud}. El \emph{estimador de máxima verosimilitud} (EMV) es la matriz estocástica que maximiza $\ell$.
\end{definicion}

El logaritmo convierte productos en sumas y evita el subdesbordamiento numérico: con $10^{3}$ factores iguales a $0{,}9$ el producto vale aproximadamente $10^{-46}$.

\section{El estimador de máxima verosimilitud}

\begin{teorema}[EMV del modelo de bigramas]\label{teo:emv}
Supóngase $N(u)>0$ para todo $u$. El estimador de máxima verosimilitud de la matriz de transición es la frecuencia relativa
\[
\widehat{P}_{uv}=\frac{N(u,v)}{N(u)}.
\]
\end{teorema}

\addtocounter{definicion}{1}%

\begin{table}[h]
\centering\small
\caption{Conteos $N(u,v)$ del \cref{ej:conteos} (fila: token anterior; columna: token siguiente).}\label{tab:conteos}
\begin{tabular}{@{}c|ccccc|c@{}}
\toprule
$u\backslash v$ & \fin & \texttt{a} & \texttt{l} & \texttt{n} & \texttt{u} & $N(u)$ \\
\midrule
\fin & 0 & 2 & 1 & 0 & 0 & 3 \\
\texttt{a} & 3 & 0 & 1 & 1 & 0 & 5 \\
\texttt{l} & 0 & 1 & 0 & 0 & 1 & 2 \\
\texttt{n} & 0 & 2 & 0 & 0 & 0 & 2 \\
\texttt{u} & 0 & 0 & 0 & 1 & 0 & 1 \\
\bottomrule
\end{tabular}
\end{table}

\section{Ceros y suavizado aditivo}

El EMV asigna probabilidad nula a todo bigrama que no aparece en el corpus. Si un texto de evaluación contiene uno de ellos, su probabilidad es~0 y la pérdida del \cref{tema:5} es infinita. Este fenómeno es una forma de sobreajuste~\cite[§3.5]{jm3}.

\begin{definicion}[Suavizado aditivo]\label{def:laplace}
Sea $\alpha>0$. El estimador con \emph{suavizado aditivo} es
\[
\widehat{P}_{\alpha}(v\mid u)=\frac{N(u,v)+\alpha}{N(u)+\alpha K}.
\]
Con $\alpha=1$ recibe el nombre de \emph{suavizado de Laplace}; con $0<\alpha<1$, de \emph{suavizado add-$k$}~\cite[§3.6.1–3.6.2]{jm3}.
\end{definicion}

\begin{proposicion}\label{prop:laplace}
Para cada $u$, $\widehat{P}_{\alpha}(\cdot\mid u)$ es una distribución de probabilidad con todas sus entradas estrictamente positivas. Además, coincide con la media a posteriori de la fila $u$ cuando se adopta una distribución a priori de Dirichlet de parámetros $(\alpha,\dots,\alpha)$.
\end{proposicion}

\addtocounter{definicion}{1}%

\section{Interpolación}

Otra estrategia consiste en combinar estimadores de distinto orden~\cite[§3.6.3]{jm3}.

\begin{definicion}[Interpolación lineal]
Dados un estimador de bigramas $\widehat{P}_2(v\mid u)$ y uno de unigramas $\widehat{P}_1(v)$, y pesos $\lambda_1,\lambda_2\ge 0$ con $\lambda_1+\lambda_2=1$, el estimador \emph{interpolado} es
\[
\widehat{P}_{\mathrm{int}}(v\mid u)=\lambda_2\,\widehat{P}_2(v\mid u)+\lambda_1\,\widehat{P}_1(v).
\]
\end{definicion}

Es una distribución, por ser combinación convexa de distribuciones. Los pesos se eligen maximizando la verosimilitud de un conjunto de validación separado del de entrenamiento.

